\chapter{Appendix: Sets, Sums, and Combinatorial Identities}

This appendix collects foundational mathematical tools used throughout the text:
\emph{sets} (the language of events), \emph{sums} (the evaluation of probabilities
and expectations), and \emph{combinatorics} (counting outcomes). Derivations in
probability frequently invoke summation formulas and combinatorial identities;
this chapter gathers their definitions, proofs, and standard applications in one place.

You do not need to read this front to back. Treat it as a self-contained reference:
each result is stated cleanly, motivated geometrically or combinatorially, and
derived explicitly.

\section{Sets: the language of events}
\label{ch39:sets}

Before probability assigns numbers to anything, it has to say \emph{what} it
is assigning numbers to. The answer is always a set. A \emph{sample space}
$\Omega$ is the set of all possible outcomes of an experiment, and an
\emph{event} is just a subset $A \subseteq \Omega$. ``The die shows an even
number'' is the set $\set{2,4,6}$. ``At least one head in three flips'' is a
subset of the eight possible flip-sequences. Once you see that events
\emph{are} sets, the algebra of probability becomes the algebra of sets, and
the standard laws of probability inherit their structure from fundamental
laws about unions and intersections.

\begin{intuitionbox}
Think of $\Omega$ as a big rectangle --- the universe of everything that
could happen. An event is a region drawn inside it. ``And'' is the overlap of
two regions, ``or'' is their combined area, and ``not'' is everything outside
a region. Probability later puts weight on these regions, but the
geometry of Boolean operations is fixed before any probabilities are assigned.
Almost every identity in this section can be visualized directly in a Venn diagram.
\end{intuitionbox}

\subsection{The basic operations}

Let $A, B \subseteq \Omega$. The fundamental set operations are:
\begin{itemize}
  \item \textbf{Union} $A \cup B = \set{\omega : \omega \in A \text{ or }
        \omega \in B}$ --- everything in at least one of them. The ``or'' is
        \emph{inclusive}: outcomes in both count too.
  \item \textbf{Intersection} $A \cap B = \set{\omega : \omega \in A \text{
        and } \omega \in B}$ --- everything in both at once.
  \item \textbf{Complement} $A^c = \set{\omega \in \Omega : \omega \notin A}$
        --- everything \emph{not} in $A$. The complement always lives inside
        the fixed universe $\Omega$.
\end{itemize}
Two more derived notions appear constantly. The \emph{difference}
$A \setminus B = A \cap B^c$ is ``in $A$ but not $B$.'' Two events are
\emph{disjoint} (or \emph{mutually exclusive}) if $A \cap B = \emptyset$ ---
they cannot both happen. A collection $A_1, A_2, \dots$ \emph{partitions}
$\Omega$ if the sets are pairwise disjoint and their union is all of $\Omega$;
this is exactly the setup behind the law of total probability.

\begin{figure}[H]
\centering
\begin{tikzpicture}[scale=1.0]
  \begin{scope}
    \clip (-1.3,-1.3) rectangle (3.0,1.6);
    \fill[sourceblue!18] (0,0) circle (1.2);
    \fill[targetgreen!22] (1.5,0) circle (1.2);
    \begin{scope}
      \clip (0,0) circle (1.2);
      \fill[structurepurple!35] (1.5,0) circle (1.2);
    \end{scope}
  \end{scope}
  \draw[thick] (-1.3,-1.3) rectangle (3.0,1.6);
  \draw[thick,sourceblue] (0,0) circle (1.2);
  \draw[thick,targetgreen] (1.5,0) circle (1.2);
  \node[sourceblue, font=\bfseries] at (-0.55,0.5) {$A$};
  \node[targetgreen, font=\bfseries] at (2.05,0.5) {$B$};
  \node[structurepurple!80!black] at (0.75,0) {\small $A\cap B$};
  \node at (2.7,1.35) {$\Omega$};
\end{tikzpicture}
\caption{Two events as overlapping regions inside the universe $\Omega$. The
purple lens is $A \cap B$; the entire shaded area is $A \cup B$; everything
outside both circles is $(A \cup B)^c = A^c \cap B^c$.}
\label{ch39:venn}
\end{figure}

\subsection{The identities you actually use}

The following laws are the grammar of set algebra. Each one has a one-line
``why'' that you can check by chasing a single outcome $\omega$ through both
sides --- this technique, \emph{element-chasing}, proves every set identity:
to show $S = T$, show that $\omega \in S$ if and only if $\omega \in T$.

\begin{formalbox}
For all events $A, B, C \subseteq \Omega$:
\begin{align*}
\textbf{Commutative:}\quad & A \cup B = B \cup A, \qquad A \cap B = B \cap A.\\
\textbf{Associative:}\quad & (A \cup B)\cup C = A \cup (B \cup C), \quad
  (A \cap B)\cap C = A \cap (B \cap C).\\
\textbf{Distributive:}\quad & A \cap (B \cup C) = (A\cap B)\cup(A\cap C),\\
  & A \cup (B \cap C) = (A\cup B)\cap(A\cup C).\\
\textbf{Identity:}\quad & A \cup \emptyset = A, \qquad A \cap \Omega = A.\\
\textbf{Complement:}\quad & A \cup A^c = \Omega, \qquad A \cap A^c = \emptyset,
  \qquad (A^c)^c = A.\\
\textbf{De Morgan:}\quad & (A \cup B)^c = A^c \cap B^c, \qquad
  (A \cap B)^c = A^c \cup B^c.
\end{align*}
\end{formalbox}

The last line is the famous pair of \emph{De Morgan's laws}, and they deserve
a moment because probability uses them constantly --- often to turn a hard
``at least one'' question into an easy ``none'' question.

\begin{whybox}
\textbf{De Morgan, in words and then in symbols.} Saying ``it is \emph{not}
the case that ($A$ or $B$)'' is exactly the same as saying ``not-$A$ \emph{and}
not-$B$.'' If neither rained-on-Monday nor rained-on-Tuesday happened, then it
failed to rain on Monday \emph{and} it failed to rain on Tuesday.

Element-chasing makes it airtight. Take $\omega \in (A\cup B)^c$. By
definition of complement, $\omega \notin A \cup B$, which means $\omega$ is in
neither $A$ nor $B$, i.e.\ $\omega \in A^c$ and $\omega \in B^c$, i.e.\
$\omega \in A^c \cap B^c$. Every step is reversible, so the two sets contain
exactly the same outcomes and are therefore equal. The general form for any
collection is
\[
  \Big(\bigcup_i A_i\Big)^c = \bigcap_i A_i^c, \qquad
  \Big(\bigcap_i A_i\Big)^c = \bigcup_i A_i^c.
\]
\end{whybox}

\begin{happeningbox}
\textbf{Why probabilists love the complement trick.} Suppose you want the
probability that \emph{at least one} of $n$ independent events occurs. Computed
directly, ``at least one'' is a sprawling union with many overlaps. But its
complement, by De Morgan, is ``none of them occur'' --- a single clean
intersection. So
\[
  \P\!\left(\bigcup_{i=1}^n A_i\right) = 1 - \P\!\left(\bigcap_{i=1}^n
  A_i^c\right),
\]
and for independent events the right-hand side factors into a product. This is
the engine behind the birthday problem (Chapter~2) and behind any
``probability that the system fails'' reliability calculation. De Morgan's set
identity reduces the event to a complement, and the probability rule $1-\P(\cdot)$ finishes the step.
\end{happeningbox}

\subsection{Cartesian products and power sets}

Two constructions show up whenever we model several things at once or talk
about ``all possible events.''

The \emph{Cartesian product} $A \times B = \set{(a,b) : a \in A,\ b \in B}$ is
the set of ordered pairs. If you roll a die and flip a coin, the joint sample
space is $\set{1,\dots,6} \times \set{H,T}$, which has $6 \times 2 = 12$
outcomes. This multiplication of sizes is the
\emph{multiplication principle} of counting, which we meet again in
Section~\ref{ch39:counting}. For finite sets, $\abs{A \times B} = \abs{A}\cdot
\abs{B}$.

The \emph{power set} $\mathcal{P}(\Omega)$ is the set of \emph{all} subsets of
$\Omega$, including $\emptyset$ and $\Omega$ itself. When $\Omega$ is finite
with $\abs{\Omega} = n$, a subset is built by an independent yes/no choice for
each element --- ``in or out'' --- so there are $2^n$ subsets:
\[
  \abs{\mathcal{P}(\Omega)} = 2^{\abs{\Omega}}.
\]
This is why a finite probability model on $n$ outcomes has $2^n$ possible
events. (For infinite $\Omega$ we cannot use \emph{all} subsets; we restrict
to a $\sigma$-algebra. That subtlety belongs to the measure-theory discussion
in the main text, but it is the reason the careful definition of an event is a
little fussier than ``any subset.'')

\begin{connectionbox}
The ``in-or-out'' bijection between subsets of an $n$-element set and binary
strings of length $n$ is the same idea that powers the binomial theorem in
Section~\ref{ch39:sums}: choosing which factors contribute is exactly choosing
a subset. Counting subsets by their size will hand us $\sum_k \binom{n}{k} =
2^n$ for free. Keep that picture in mind --- sets and sums are two views of
one object.
\end{connectionbox}

\section{Sigma-notation and the sums we cannot live without}
\label{ch39:sums}

\subsection{Reading and bending $\sum$}

The symbol $\sum_{k=m}^{n} a_k$ means ``add up $a_k$ as the index $k$ runs
from $m$ to $n$.'' The index is a \emph{dummy}: $\sum_{k=1}^n a_k$ and
$\sum_{j=1}^n a_j$ are the same number. Three manipulations cover nearly every
move you will ever make.

\begin{itemize}
  \item \textbf{Linearity.} $\displaystyle \sum_k (c\,a_k + b_k) = c\sum_k a_k
        + \sum_k b_k$. Constants pull out; sums split apart. This is exactly
        why expectation is linear --- $\E$ is a weighted sum.
  \item \textbf{Reindexing (shifting).} Let $j = k - m$ to slide the limits:
        $\sum_{k=m}^{n} a_k = \sum_{j=0}^{n-m} a_{j+m}$. Useful for lining up
        two sums so their terms match.
  \item \textbf{Swapping double sums.} $\sum_{i}\sum_{j} a_{ij} =
        \sum_{j}\sum_{i} a_{ij}$ when the terms are nonnegative or absolutely
        summable. This is the discrete twin of Fubini's theorem and underlies
        ``$\E[\sum X_i] = \sum \E[X_i]$'' even when the $X_i$ are dependent.
\end{itemize}

\begin{trapbox}
\textbf{Index bounds and off-by-one errors.} The sum
$\sum_{k=0}^{n} a_k$ has $n+1$ terms, whereas $\sum_{k=1}^{n} a_k$ has $n$ terms.
When reindexing, change both the limits and the summand together, and verify
by writing out the first and last terms explicitly.
\end{trapbox}

\subsection{The arithmetic series}

Summing the first $n$ positive integers yields a foundational formula:

\begin{proposition}[Gauss's sum]
\label{ch39:gauss}
For every $n \ge 1$,
\[
  \sum_{k=1}^{n} k = 1 + 2 + \cdots + n = \frac{n(n+1)}{2}.
\]
\end{proposition}

\begin{whybox}
Write the sum forwards and backwards and add columnwise:
\[
\begin{array}{ccccccc}
S &=& 1 &+& 2 &+ \cdots +& n,\\
S &=& n &+& (n-1) &+ \cdots +& 1.
\end{array}
\]
Each of the $n$ columns adds to $n+1$, so $2S = n(n+1)$, giving $S =
n(n+1)/2$. There is no induction needed --- just a
reflection. A famous legend tells of young Carl Friedrich Gauss pairing $1{+}100,
2{+}99, \dots$ into fifty pairs of $101$ to sum the first hundred integers in
seconds; whether historical fact or charming folklore, this pairing principle is
the exact geometric reflection at work in the formula. A standard telescoping
argument gives $\sum_{k=1}^n k^2 = n(n+1)(2n+1)/6$, which you meet whenever you
compute the variance of a discrete uniform distribution.
\end{whybox}

\subsection{The geometric series}
\label{ch39:geom-series}

The geometric series is central to probability: it governs the geometric distribution,
expected waiting times, renewal processes, and stationary distributions of Markov chains.

\begin{proposition}[Finite and infinite geometric series]
\label{ch39:geom}
For $r \neq 1$,
\[
  \sum_{k=0}^{n-1} r^k = \frac{1 - r^n}{1 - r}.
\]
If $\abs{r} < 1$, then $r^n \to 0$ and the infinite series converges:
\[
  \sum_{k=0}^{\infty} r^k = \frac{1}{1 - r}.
\]
\end{proposition}

\begin{whybox}
\textbf{The telescoping trick.} Let $S = 1 + r + r^2 + \cdots + r^{n-1}$.
Multiply by $r$: $rS = r + r^2 + \cdots + r^n$. Subtract:
\[
  S - rS = 1 - r^n \quad\Longrightarrow\quad S(1-r) = 1 - r^n
  \quad\Longrightarrow\quad S = \frac{1-r^n}{1-r}.
\]
Almost everything cancels --- the middle terms \emph{telescope}, leaving only
the first term of $S$ and the last term of $rS$. For $\abs{r}<1$ the tail
$r^n$ shrinks to nothing, so letting $n \to \infty$ leaves $1/(1-r)$. The
condition $\abs{r}<1$ is not a technicality you may ignore: if $\abs{r}\ge 1$
the terms do not shrink and the sum is infinite or undefined.
\end{whybox}

A relative that appears whenever we compute the \emph{mean} of a geometric
random variable is the differentiated series. Differentiating $\sum r^k =
1/(1-r)$ term by term with respect to $r$ (legal inside the radius of
convergence) gives
\[
  \sum_{k=1}^{\infty} k\, r^{k-1} = \frac{1}{(1-r)^2}, \qquad\text{hence}\qquad
  \sum_{k=1}^{\infty} k\, r^{k} = \frac{r}{(1-r)^2} \quad (\abs{r}<1).
\]

\begin{recoverybox}
\textbf{Pause and think:} Why does $\sum_{k=1}^\infty k\,r^k$ start at $k=1$ while $\sum_{k=0}^\infty r^k$ starts at $k=0$? At $k=0$, the term $k\,r^k = 0\cdot r^0 = 0$ contributes nothing. If $r = 1/2$, what is $\sum_{k=1}^\infty k\,(1/2)^k$? Plugging into $r/(1-r)^2$ gives $(1/2)/(1/4) = 2$ --- exactly the mean of a $\Geom(1/2)$ distribution!
\end{recoverybox}

\begin{figure}[H]
\centering
\begin{tikzpicture}
  \begin{axis}[
      width=11.5cm, height=6.2cm,
      title={Partial sums of the geometric series climbing to limit $1/(1-r)$},
      title style={font=\small},
      xlabel={number of terms $n$},
      ylabel={partial sum $\sum_{k=0}^{n-1} r^k$},
      xmin=0, xmax=12, ymin=0, ymax=2.40,
      grid=both, grid style={gray!18},
      samples at={0,1,...,12},
  ]
    \addplot[ycomb, mark=*, sourceblue, thick] {(1-0.5^x)/(1-0.5)};
    \addplot[densityred, dashed, thick, domain=0:12]{2};
    \node[densityred, anchor=south east] at (axis cs:11.8, 2.03) {\small limit $1/(1-r)=2$};
  \end{axis}
\end{tikzpicture}
\caption{Partial sums of the geometric series with $r = 0.5$ marching up to
their limit $1/(1-r) = 2$. The gap remaining after $n$ terms is exactly
$r^n/(1-r)$, which halves at every step.}
\label{ch39:geomfig}
\end{figure}

\begin{applicationbox}
\textbf{Where does $1/(1-r)$ become a probability statement?} Let $X$ be the
number of independent trials, each succeeding with probability $p$, up to and
including the first success (a geometric random variable). Then $\P(X = k) =
(1-p)^{k-1}p$ for $k = 1,2,\dots$. For this to be a valid distribution the
probabilities must sum to $1$, and indeed, with $r = 1-p$,
\[
  \sum_{k=1}^{\infty}(1-p)^{k-1}p = p\sum_{j=0}^{\infty} r^{j}
  = p\cdot\frac{1}{1-r} = \frac{p}{p} = 1.
\]
The geometric series is what \emph{certifies} the geometric distribution. And
the differentiated series gives its mean: $\E[X] = \sum_k k(1-p)^{k-1}p =
p/(1-r)^2 = 1/p$. So if a success has probability $1/6$, you wait $6$ trials on
average --- a fact you can now derive rather than recall.
\end{applicationbox}

\begin{figure}[H]
\centering
\begin{tikzpicture}
  \begin{axis}[
      width=11.5cm, height=5.6cm,
      title={The geometric distribution pmf $\P(X=k) = (1-p)^{k-1}p$},
      title style={font=\small},
      ybar, bar width=10pt,
      xlabel={trial of first success $k$},
      ylabel={$\P(X=k)=(1-p)^{k-1}p$},
      xmin=0.4, xmax=10.6, ymin=0, ymax=0.38,
      xtick={1,2,...,10},
      grid=both, grid style={gray!18},
      legend pos=north east,
  ]
    \addplot[draw=sourceblue!70!black, fill=sourceblue!35, area legend]
      coordinates {
        (1,0.33333)(2,0.22222)(3,0.14815)(4,0.09877)(5,0.06584)
        (6,0.04390)(7,0.02926)(8,0.01951)(9,0.01301)(10,0.00867)};
    \addlegendentry{$p=1/3$}
  \end{axis}
\end{tikzpicture}
\caption{The geometric pmf with $p=1/3$. The bar heights are exactly the terms
of the geometric series $p\sum_{k\ge 1} r^{k-1}$ with $r=1-p=2/3$; each bar is
$r$ times the one before it, and the whole row of bars sums to $1$. The mean
$\E[X]=1/p=3$ sits to the right of the tall first bar because the long thin
tail pulls the balance point outward.}
\label{ch39:geompmf}
\end{figure}

\subsection{The exponential series}

The number $e$ and its series are the bridge from discrete to continuous
probability --- they convert the binomial into the Poisson and define the
exponential and normal densities.

\begin{proposition}[Exponential series]
\label{ch39:exp}
For every real (indeed complex) $x$,
\[
  e^{x} = \sum_{k=0}^{\infty} \frac{x^{k}}{k!}
  = 1 + x + \frac{x^2}{2} + \frac{x^3}{6} + \cdots,
\]
and equivalently $\displaystyle e^{x} = \lim_{n\to\infty}\Big(1 +
\frac{x}{n}\Big)^{n}$.
\end{proposition}

\begin{happeningbox}
\textbf{Why two formulas for the same thing, and why probability needs both.}
The series $\sum x^k/k!$ is the Taylor expansion of $e^x$; it is what you
differentiate and integrate. The limit $(1+x/n)^n$ is the
``compound-interest'' face of $e$: split a rate $x$ into $n$ tiny increments
and compound. Probability uses the second form to derive the Poisson
distribution as the limit of a binomial: take $n$ trials each with tiny
success probability $\lambda/n$, and watch
\[
  \Big(1 - \frac{\lambda}{n}\Big)^{n} \longrightarrow e^{-\lambda}.
\]
That single limit explains why rare events --- radioactive decays, typos
per page, mutations per genome --- follow a Poisson law. This limit connects
discrete counting processes with continuous exponential models.
\end{happeningbox}

\begin{connectionbox}
Setting $x = -\lambda$ shows $\sum_{k=0}^\infty e^{-\lambda}\lambda^k/k! =
e^{-\lambda}e^{\lambda} = 1$, which is exactly the statement that the Poisson
pmf sums to one. As with the geometric distribution, a series identity is what
makes a distribution legal. The calculus you need to differentiate these
series term by term --- power series, radius of convergence, term-by-term
differentiation --- is developed in the companion calculus volume; here we use
it as a trusted tool.
\end{connectionbox}

\subsection{The binomial theorem}

This is the algebraic identity that ties counting to sums. It is the reason
the binomial \emph{distribution} is named after it.

\begin{theorem}[Binomial theorem]
\label{ch39:binom}
For any numbers $x, y$ and any integer $n \ge 0$,
\[
  (x + y)^{n} = \sum_{k=0}^{n} \binom{n}{k} x^{k} y^{\,n-k},
\]
where $\binom{n}{k} = \dfrac{n!}{k!\,(n-k)!}$ is the number of $k$-element
subsets of an $n$-element set.
\end{theorem}

\begin{whybox}
\textbf{Combinatorial expansion.} Write $(x+y)^n = (x+y)(x+y)\cdots(x+y)$ with
$n$ factors. To expand, we select either $x$ or $y$ from \emph{each} factor
and multiply the choices; the full product is the sum over all $2^n$ such
choices. A term equals $x^k y^{n-k}$ precisely when we choose $x$ from $k$ of
the factors and $y$ from the other $n-k$. The number of ways to choose
\emph{which} $k$ factors give $x$ is $\binom{n}{k}$. Collecting like terms
yields the coefficient $\binom{n}{k}$ in front of $x^k y^{n-k}$, directly
counting the number of ways that term is formed.
\end{whybox}

Two specializations are worth keeping at your fingertips. Setting $x = y = 1$:
\[
  2^{n} = \sum_{k=0}^{n}\binom{n}{k},
\]
the count of all subsets, matching the power-set count from
Section~\ref{ch39:sets}. Setting $x = 1,\ y = -1$ (for $n \ge 1$):
\[
  0 = \sum_{k=0}^{n}(-1)^{k}\binom{n}{k},
\]
so the alternating sum of a row of Pascal's triangle vanishes --- the even-
and odd-sized subsets exactly balance.

\begin{applicationbox}
\textbf{The binomial distribution sums to one.} If $X \sim \Bin(n,p)$ counts
successes in $n$ independent trials, then $\P(X=k) = \binom{n}{k}p^k(1-p)^{n-k}$.
Summing over $k$ and applying the binomial theorem with $x=p$, $y=1-p$:
\[
  \sum_{k=0}^{n}\binom{n}{k}p^{k}(1-p)^{n-k} = \big(p + (1-p)\big)^{n}
  = 1^{n} = 1.
\]
Once again an algebraic identity guarantees a probability distribution is
well-defined. The pattern across this whole section is unmistakable: every
named distribution is licensed by a named sum.
\end{applicationbox}

\section{Combinatorial identities: counting the same thing two ways}
\label{ch39:counting}

\subsection{Fundamental counting principles}

Combinatorial enumeration relies on two basic principles:

\textbf{Multiplication principle.} If a process has two stages, with $m$ ways
to do the first and (for each of those) $n$ ways to do the second, there are
$m \cdot n$ ways total. Generalizes to any number of stages. This is why an
$n$-element set has $n!$ orderings: $n$ choices for first place, $n-1$ for
second, and so on.

\textbf{Addition principle.} If the outcomes split into disjoint cases, the
total count is the sum of the case counts. This is just $\abs{A \cup B} =
\abs{A} + \abs{B}$ for disjoint $A, B$ --- the set law from
Section~\ref{ch39:sets} reappearing as a counting rule.

\begin{figure}[H]
\centering
\begin{tikzpicture}[
    level 1/.style={sibling distance=34mm, level distance=12mm},
    level 2/.style={sibling distance=12mm, level distance=12mm},
    every node/.style={font=\small},
    edge from parent/.style={draw, gray!70, -{Latex[length=1.6mm]}},
  ]
  \node[draw, circle, fill=sourceblue!10, inner sep=1.5pt] {start}
    child {node[draw, circle, fill=sourceblue!18, inner sep=1.5pt] {$1$}
      child {node[draw, fill=targetgreen!18, inner sep=2pt] {$H$}}
      child {node[draw, fill=targetgreen!18, inner sep=2pt] {$T$}}}
    child {node[draw, circle, fill=sourceblue!18, inner sep=1.5pt] {$2$}
      child {node[draw, fill=targetgreen!18, inner sep=2pt] {$H$}}
      child {node[draw, fill=targetgreen!18, inner sep=2pt] {$T$}}}
    child {node[draw, circle, fill=sourceblue!18, inner sep=1.5pt] {$3$}
      child {node[draw, fill=targetgreen!18, inner sep=2pt] {$H$}}
      child {node[draw, fill=targetgreen!18, inner sep=2pt] {$T$}}};
\end{tikzpicture}
\caption{The multiplication principle as a tree. Rolling a $3$-sided die (blue,
$3$ choices) and then flipping a coin (green, $2$ choices) yields
$3 \times 2 = 6$ leaves, one per outcome. Each layer of branching multiplies
the count; this is the same product that makes $\abs{A\times B}=\abs{A}\,
\abs{B}$ in Section~\ref{ch39:sets}.}
\label{ch39:treefig}
\end{figure}

From these we get the four basic counts. Choosing $k$ items from $n$:

\begin{center}
\begin{tabular}{lcc}
\toprule
 & \textbf{order matters} & \textbf{order ignored}\\
\midrule
\textbf{no repetition} & $\dfrac{n!}{(n-k)!}$ & $\dbinom{n}{k} =
  \dfrac{n!}{k!(n-k)!}$\\[1.2em]
\textbf{with repetition} & $n^{k}$ & $\dbinom{n+k-1}{k}$\\
\bottomrule
\end{tabular}
\end{center}

\begin{intuitionbox}
The binomial coefficient $\binom{n}{k}$ --- read ``$n$ choose $k$'' ---
counts the $k$-element subsets of an $n$-element set. To see why $\binom{n}{k}
= \frac{n!}{k!(n-k)!}$, consider ordering all $n$ items ($n!$ ways). The
internal ordering of the $k$ chosen items does not change the subset ($k!$
orderings), and the internal ordering of the $n-k$ unchosen items is likewise
irrelevant ($(n-k)!$ orderings). Dividing by these redundant orderings leaves
the subset count.
\end{intuitionbox}

\subsection{Identities and how to read them}

A standard technique for proving combinatorial identities is a
\emph{combinatorial proof}: find a set of objects, count it two different
ways, and equate the resulting expressions. The following identities appear
frequently in probability calculations, each accompanied by its two-count
interpretation.

\begin{formalbox}
For integers $0 \le k \le n$:
\begin{align}
\binom{n}{k} &= \binom{n}{n-k}
  &&\text{(symmetry)} \label{ch39:sym}\\[0.2em]
\binom{n}{k} &= \binom{n-1}{k-1} + \binom{n-1}{k}
  &&\text{(Pascal's rule)} \label{ch39:pascal}\\[0.2em]
\sum_{k=0}^{n}\binom{n}{k} &= 2^{n}
  &&\text{(row sum)} \label{ch39:rowsum}\\[0.2em]
\sum_{k=0}^{r}\binom{m}{k}\binom{n}{r-k} &= \binom{m+n}{r}
  &&\text{(Vandermonde)} \label{ch39:vander}\\[0.2em]
k\binom{n}{k} &= n\binom{n-1}{k-1}
  &&\text{(absorption / committee--chair)} \label{ch39:absorb}
\end{align}
\end{formalbox}

\begin{whybox}
\textbf{Symmetry \eqref{ch39:sym}.} Choosing which $k$ people to \emph{invite}
is the same as choosing which $n-k$ to \emph{leave out}. Each choice of an
invited group corresponds to exactly one excluded group, so the two counts
must match. (Algebraically, swap $k$ and $n-k$ in the factorial formula and
nothing changes.)

\textbf{Pascal's rule \eqref{ch39:pascal}.} Count the $k$-subsets of
$\set{1,\dots,n}$ by whether they contain element $n$. If they \emph{do}, the
remaining $k-1$ members come from the other $n-1$ elements:
$\binom{n-1}{k-1}$ ways. If they \emph{do not}, all $k$ come from the other
$n-1$: $\binom{n-1}{k}$ ways. The two cases are disjoint and exhaustive, so by
the addition principle they sum to $\binom{n}{k}$. This single recurrence
\emph{is} Pascal's triangle.

\textbf{Vandermonde's identity \eqref{ch39:vander}.} You have $m$ women and
$n$ men and want to form a committee of $r$ people: $\binom{m+n}{r}$ ways.
Alternatively, decide first that the committee has $k$ women (chosen
$\binom{m}{k}$ ways) and $r-k$ men ($\binom{n}{r-k}$ ways), then sum over all
possible $k$. Same committees, counted two ways.

\textbf{Absorption \eqref{ch39:absorb}.} Count ways to pick a $k$-person
committee \emph{and} designate one of its members as chair. Pick the committee,
then the chair: $\binom{n}{k}\cdot k$. Or pick the chair from all $n$ people
first, then fill the remaining $k-1$ seats from the other $n-1$:
$n\binom{n-1}{k-1}$. Equal counts. This identity is exactly the algebraic step
that produces the mean of the binomial distribution, $\E[X] = np$.
\end{whybox}

\begin{figure}[H]
\centering
\begin{tikzpicture}[x=0.95cm,y=0.9cm,font=\small]
  % Pascal's triangle, rows 0..5
  \foreach \r [count=\ri from 0] in {
      {1},
      {1,1},
      {1,2,1},
      {1,3,3,1},
      {1,4,6,4,1},
      {1,5,10,10,5,1}} {
    \foreach \v [count=\ci from 0] in \r {
      \node[draw=sourceblue!55, fill=sourceblue!8, rounded corners=2pt,
            minimum size=0.62cm, inner sep=1pt]
        at ({\ci - \ri/2},{-\ri}) {\v};
    }
  }
  % highlight Pascal's rule: 3+3 -> 6
  \draw[densityred,thick,-{Latex[length=1.5mm]}] (-0.35,-3.35) -- (-0.12,-3.65);
  \draw[densityred,thick,-{Latex[length=1.5mm]}] (0.35,-3.35) -- (0.12,-3.65);
  \node[densityred,anchor=west] at (2.6,-4.0) {\small $3+3=6$};
\end{tikzpicture}
\caption{Pascal's triangle. Each entry is the sum of the two above it
(Pascal's rule, \eqref{ch39:pascal}); row $n$ holds $\binom{n}{0},\dots,
\binom{n}{n}$ and sums to $2^n$.}
\label{ch39:pascalfig}
\end{figure}

\begin{paradoxbox}
\textbf{The birthday ``paradox'' is pure counting.} In a room of just $23$
people, the probability that two share a birthday already exceeds one half ---
shocking, because $23$ feels tiny against $365$ days. The resolution lives
entirely in this section. Use the complement trick (Section~\ref{ch39:sets}):
the chance that all $23$ birthdays are \emph{distinct} is, by the
multiplication principle,
\[
  \frac{365}{365}\cdot\frac{364}{365}\cdots\frac{343}{365}
  = \frac{365!/(365-23)!}{365^{23}} \approx 0.493,
\]
so a match has probability $\approx 0.507$. The trap is that we instinctively
compare ourselves to everyone else, but the relevant quantity is the number of
\emph{pairs}, $\binom{23}{2} = 253$ --- and $253$ chances at a $1/365$ event
is no longer rare. Nothing here is paradoxical; it is the gap between linear
intuition (people) and quadratic reality (pairs).
\end{paradoxbox}

\subsection{Inclusion--exclusion}

When events overlap, you cannot simply add their sizes --- you would
double-count the overlap. Fixing this systematically gives one of the most
useful counting (and probability) formulas in the book.

\begin{theorem}[Inclusion--exclusion]
\label{ch39:incexc}
For any events $A_1,\dots,A_n$,
\[
  \abs[\Big]{\bigcup_{i=1}^n A_i}
  = \sum_i \abs{A_i} - \sum_{i<j}\abs{A_i \cap A_j}
  + \sum_{i<j<k}\abs{A_i \cap A_j \cap A_k} - \cdots
  + (-1)^{n+1}\abs[\Big]{\bigcap_{i=1}^n A_i}.
\]
The identical statement holds with $\abs{\cdot}$ replaced by $\P(\cdot)$.
\end{theorem}

\begin{whybox}
For two sets it is the picture in Figure~\ref{ch39:venn}: $\abs{A \cup B} =
\abs{A} + \abs{B} - \abs{A \cap B}$, because adding the two circles counts the
lens twice, so we subtract it once. For three sets you add back the triple
overlap you over-subtracted, and so on --- each correction fixes the previous
over- or under-count. The clean proof tracks a single element belonging to
exactly $m$ of the sets and checks that the alternating sum counts it
\[
  \binom{m}{1} - \binom{m}{2} + \binom{m}{3} - \cdots = 1 - (1-1)^m = 1
\]
time (using the binomial theorem with $x=-1,y=1$!), exactly as it should. The
combinatorial identity \eqref{ch39:rowsum}'s alternating cousin is doing the
bookkeeping.
\end{whybox}

\begin{applicationbox}
\textbf{Inclusion--exclusion in the data center.} A search or analytics engine
constantly needs the size of a union: ``how many distinct users visited page
$A$ \emph{or} page $B$ \emph{or} page $C$?'' Storing every user id is
expensive, so systems keep a tiny probabilistic sketch (a Bloom filter or a
HyperLogLog counter) of each set. To estimate the union from the individual
sketches, engineers apply exactly Theorem~\ref{ch39:incexc}: add the
single-set cardinalities, subtract the pairwise intersections, add back the
triples. The same alternating signs that fix double-counting on a Venn diagram
fix it at the scale of billions of rows.
\end{applicationbox}

\section{Asymptotic approximations: Stirling's formula}
\label{ch39:asymptotics}

Factorials grow rapidly, making exact evaluation impractical for large $n$. In
probability, we frequently need the asymptotic behavior of $n!$ as $n \to
\infty$. Stirling's formula provides a precise approximation.

\begin{theorem}[Stirling's approximation]
\label{ch39:stirling}
As $n \to \infty$,
\[
  n! \sim \sqrt{2\pi n}\,\Big(\frac{n}{e}\Big)^{n},
\]
meaning the ratio of the two sides tends to $1$. A more precise version with
error bounds is
\[
  n! = \sqrt{2\pi n}\,\Big(\frac{n}{e}\Big)^{n}
  \exp\!\Big(\frac{1}{12n} + O(n^{-3})\Big).
\]
\end{theorem}

\begin{happeningbox}
\textbf{Where the pieces come from.} Take logs: $\ln n! = \sum_{k=1}^n \ln k$.
A sum of a slowly varying function is well approximated by an integral, and
$\int_1^n \ln x\,\dd x = n\ln n - n + 1$. That already explains the dominant
factor $(n/e)^n = e^{n\ln n - n}$: the $n \ln n$ is the area under $\ln x$, and
the $-n$ is the $-n$ in the integral. The leftover $\sqrt{2\pi n}$ comes from
the correction to replacing a sum by an integral (the trapezoid/Euler--Maclaurin
correction), and the $2\pi$ is the same $2\pi$ that appears in the normal
density, as we will see when we derive the Gaussian approximation to the binomial.
\end{happeningbox}

\begin{figure}[H]
\centering
\begin{tikzpicture}
  \begin{axis}[
      width=11.5cm, height=6.2cm,
      title={Stirling's approximation $\ln(n!) \approx n\ln n - n + (1/2)\ln(2\pi n)$},
      title style={font=\small},
      xlabel={$n$}, ylabel={$\ln(n!)$ and Stirling estimate},
      xmin=1, xmax=12, ymin=0, ymax=20,
      legend pos=north west, legend cell align=left,
      grid=both, grid style={gray!18},
  ]
    % exact ln(n!) at integers
    \addplot[only marks, mark=*, sourceblue, mark size=2pt]
      coordinates {
        (1,0)(2,0.6931)(3,1.7918)(4,3.1781)(5,4.7875)
        (6,6.5793)(7,8.5252)(8,10.6046)(9,12.8018)
        (10,15.1044)(11,17.5023)(12,19.9872)};
    % Stirling: 0.5 ln(2 pi n) + n ln n - n
    \addplot[densityred, thick, domain=1:12, samples=120]
      {0.5*ln(2*pi*x) + x*ln(x) - x};
    \addlegendentry{exact $\ln(n!)$}
    \addlegendentry{Stirling estimate}
  \end{axis}
\end{tikzpicture}
\caption{Stirling's formula tracks $\ln(n!)$ closely even for small
$n$. At $n=10$ the relative error in $n!$ itself is already under one percent.}
\label{ch39:stirlingfig}
\end{figure}

\begin{applicationbox}
\textbf{Stirling builds the bell curve.} The central limit theorem's bell
shape can be seen directly by applying Stirling to the binomial coefficient
$\binom{2m}{m}$, the central term of the symmetric binomial. Stirling gives
\[
  \binom{2m}{m} = \frac{(2m)!}{(m!)^2} \sim \frac{4^{m}}{\sqrt{\pi m}}.
\]
Normalizing the binomial $\Bin(n,\tfrac12)$ and letting $n$ grow, those
$\sqrt{\cdot}$ factors assemble into $\frac{1}{\sqrt{2\pi}}e^{-x^2/2}$ --- the
standard normal density. The $2\pi$ in Stirling's formula is the same $2\pi$
appearing in the Gaussian normalization constant.
\end{applicationbox}

\begin{historybox}
The approximation is named for Scottish mathematician James Stirling, who
published it in his 1730 treatise \emph{Methodus Differentialis}. The
impetus came from Abraham de Moivre, who needed to estimate the central
binomial coefficient $\binom{2m}{m}(1/2)^{2m}$ for large numbers of coin
flips. By approximating the sum of logarithms $\sum_{k=1}^n \ln k$ with the
integral of $\ln x$, de Moivre showed that $n! \sim C\,n^{n+1/2}e^{-n}$. He could
not identify the constant $C$, but Stirling, using John Wallis's 1655 product
for $\pi$, proved that $C = \sqrt{2\pi}$. In his 1733 supplement to \emph{The
Doctrine of Chances}, de Moivre used this estimate to prove the first version of
the Central Limit Theorem, establishing the normal curve decades before Gauss.
\end{historybox}

\begin{aestheticbox}
A unified structure connects these tools: set identities simplify event
descriptions, series identities verify probability normalizations and yield
moments, combinatorial identities enumerate outcomes, and asymptotic
approximations such as Stirling's formula connect discrete distributions to
continuous limits.
\end{aestheticbox}

These tools cover the discrete and combinatorial side: sets of outcomes,
sums of probabilities, and counting. For continuous distributions, joint
densities, and linear algebra, see Appendix~\ref{chAcla:calc-intro}.

\section{Exercises}
\label{ch39:exercises}

\begin{exercise}[Quick mental checks]
\label{ch39:ex-warmup}
Warm up with three quick mental checks:
\begin{enumerate}
  \item If $A \subseteq B$, what are $A \cap B$ and $A \cup B$? What is $A \cap B^c$?
  \item If an infinite geometric series has ratio $r = 1/3$, what is $\sum_{k=0}^\infty (1/3)^k$? What is $\sum_{k=1}^\infty (1/3)^k$?
  \item How many ways can you choose a committee of $3$ people from a group of $5$? How many subsets of all sizes does a $5$-element set have?
\end{enumerate}
\end{exercise}

\begin{exercise}[Element-chasing]
Prove the distributive law $A \cap (B \cup C) = (A \cap B) \cup (A \cap C)$ by
showing each side is a subset of the other. Then state the dual law obtained
by swapping $\cup$ and $\cap$, and explain why De Morgan's laws make the two
duals equivalent.
\end{exercise}

\begin{exercise}[Complement trick]
A machine has $5$ components that fail independently, each with probability
$0.1$. The machine works only if \emph{all} components work. Using De Morgan
and the complement, write the probability that the machine fails as
$1 - \P(\text{none fail})$ and compute it. Why is this far easier than summing
over the ways at least one component fails?
\end{exercise}

\begin{exercise}[Geometric series]
Sum each series, stating any convergence condition you use:
(a) $\sum_{k=0}^{\infty}(2/3)^k$;\;
(b) $\sum_{k=2}^{\infty}(2/3)^k$;\;
(c) $\sum_{k=1}^{\infty}k(2/3)^{k}$.
For (c), identify which derived identity from Section~\ref{ch39:sums} you are
using.
\end{exercise}

\begin{exercise}[Distributions from sums]
Verify by a series identity that each is a valid probability distribution
(probabilities are nonnegative and sum to $1$):
(a) Poisson, $\P(X=k)=e^{-\lambda}\lambda^k/k!$ for $k=0,1,2,\dots$;\;
(b) Geometric, $\P(X=k)=(1-p)^{k-1}p$ for $k=1,2,\dots$. Name the identity
used in each case.
\end{exercise}

\begin{exercise}[Binomial theorem in action]
Use the binomial theorem to evaluate $\sum_{k=0}^{n}\binom{n}{k}2^{k}$ in
closed form. Then give a one-sentence combinatorial interpretation of your
answer in terms of choosing a subset and then doing something extra to its
elements.
\end{exercise}

\begin{exercise}[Combinatorial proof]
Give a committee-counting (two-ways) proof of Vandermonde's identity
$\sum_{k=0}^{r}\binom{m}{k}\binom{n}{r-k}=\binom{m+n}{r}$. As a special case,
set $m=n$ and $r=n$ to deduce $\sum_{k=0}^{n}\binom{n}{k}^2 = \binom{2n}{n}$,
and explain in words what is being counted.
\end{exercise}

\begin{exercise}[Pascal and the triangle]
Prove Pascal's rule $\binom{n}{k}=\binom{n-1}{k-1}+\binom{n-1}{k}$ both
(a) combinatorially (condition on whether a fixed element is chosen) and
(b) algebraically from the factorial formula. Compare the conceptual insight
of the combinatorial proof with the algebraic verification.
\end{exercise}

\begin{exercise}[Absorption and the binomial mean]
Starting from $\E[X]=\sum_{k=0}^{n}k\binom{n}{k}p^k(1-p)^{n-k}$ for
$X\sim\Bin(n,p)$, apply the absorption identity \eqref{ch39:absorb} to rewrite
$k\binom{n}{k}$, reindex the sum, and conclude $\E[X]=np$. State clearly where
the binomial theorem reappears to collapse the leftover sum.
\end{exercise}

\begin{exercise}[Inclusion--exclusion]
At a party, $70\%$ like tea, $80\%$ like coffee, and $60\%$ like both. What
fraction like at least one of the two? Now suppose a third drink is added with
its own overlaps; write the general three-set inclusion--exclusion formula you
would use, and explain why the signs alternate.
\end{exercise}

\begin{exercise}[Stirling]
(a) Use Stirling's approximation to show $\binom{2m}{m}\sim 4^m/\sqrt{\pi m}$.
(b) Estimate $\binom{100}{50}$ to one significant figure and compare the
\emph{exponent of growth} ($4^{50}$) with the polynomial correction
($1/\sqrt{50\pi}$). What does the relative smallness of the correction tell
you about why the central binomial probability still shrinks as $m$ grows?
\end{exercise}

\begin{exercise}[Conceptual synthesis]
This appendix repeatedly used algebraic and combinatorial identities to verify
probability properties. Pick any \emph{three} distributions from the main
text and, for each, name the single sum or counting identity that proves its
probabilities total $1$. Write two sentences explaining how the condition
$\sum_k p_k = 1$ connects each discrete distribution to a corresponding
algebraic series identity.
\end{exercise}
